% sec-07: Document 07, the federal compliance guide for the LLM and robotic
% workflow.  Operative for master prompt clause C.
% STAGE 3.  The fourteen-row correction map moved to mvltable, because it is
% taller than one page once its Requirement column wraps.  Two widths retuned.
\docpart{FDA Compliance Guide}{FDA LLM and Robotic Workflow National Compliance
Guide for a Phase 1 PDAC Clinical Trial Site}

\section{Purpose and Scope}

This document is a map from one clinical trial site to the federal requirements
that already govern it. It is not a submission, not a guidance document, not an
agency position, and not a claim that any workflow described here has been
reviewed or accepted. The honest form of the request that a large language model
and robotic workflow be acceptable to the agency is a map that a reviewer can
check, requirement by requirement, against artifacts that exist.

Its scope is the site described in the front matter of this package, operating
the Phase 1 protocol at document 13 §2. Nothing here applies to a marketed
device, to a commercial software product, or to any use of a model outside an
investigational setting.

\section{The Regulatory Position of the Site}

The trial is a combined investigational new drug and investigational device
investigation \cite{phase1protocol}. The filing is the 172-page investigational
new drug application deposited in July 2026, which follows the section skeleton
established by the adapted Part 312 framework \cite{indapp,cfr312adapt}.

\begin{mvtable}
\begin{tabularx}{\textwidth}{>{\raggedright\arraybackslash}p{4.2cm} >{\raggedright\arraybackslash}p{3.4cm} >{\raggedright\arraybackslash}X}
\headrow \thead{Component} & \thead{Pathway} & \thead{Why it sits there}\\
\midrule
Perioperative daraxonrasib & Investigational new drug, 21 CFR 312 & An investigational agent administered under a protocol \cite{cfr312} \\
Robotic surgical platform & Investigational device exemption, 21 CFR 812, or a cleared device used within its labeling & Determined at the site agreement; the configuration is not specified before then \\
LLM advisory output & Neither, while it remains advisory-only & It produces no output that reaches a participant except through an affirmative human act \\
Electronic records and signatures & 21 CFR 11 & Prompts, completions and dispositions are source records \cite{cfr11} \\
Financial disclosure & 21 CFR 54 & Decides the sponsor-investigator firewall in document 14 §3 \cite{cfr54} \\
Human subject protection & 21 CFR 50 and 45 CFR 46 & Consent and protection, as adapted \cite{cfr50,cfr45part46,cfr50adapt} \\
\end{tabularx}
\end{mvtable}
\tabcapl{tab:fdaposition}{Six components and the pathway each sits on. The third
row is the one that must be defended: the advisory-only boundary in §5 is what
keeps the model out of the device pathway, and if the boundary moves, the row
moves with it.}

\section{Correction Map to 21 CFR Parts 11, 50, 54, 312, and 812}

\begin{mvltable}
\begin{xltabular}{\textwidth}{>{\raggedright\arraybackslash}p{2.6cm} >{\raggedright\arraybackslash}p{4.2cm} >{\raggedright\arraybackslash}X}
\headrow \thead{Citation} & \thead{Requirement} & \thead{Where this site discharges it}\\
\midrule
\endfirsthead
\headrow \thead{Citation} & \thead{Requirement} & \thead{Where this site discharges it}\\
\midrule
\endhead
21 CFR 11.10(a) & Validation of systems to ensure accuracy and consistency & Document 12 §6 quality management; the model version register at document 03 §3 \\
21 CFR 11.10(e) & Secure, computer-generated, time-stamped audit trails & Document 03 §5, hash-chained and verified monthly \\
21 CFR 11.10(d) & Limiting system access to authorized individuals & Document 09 §2 access control; the machine room is a restricted zone \\
21 CFR 50.20 & General requirements for informed consent & Document 02 §4 and document 13 §3 \\
21 CFR 50.25 & Elements of informed consent & Document 02 §4, including the model disclosure and the human-only pathway \\
21 CFR 54.2 & Definitions triggering financial disclosure & Document 14 §3, which assigns the investigator role wholly to the site \\
21 CFR 54.4 & Certification and disclosure statements & Document 15 §2 stewardship controls \\
21 CFR 312.23 & Investigational new drug content and format & The filed application \cite{indapp} \\
21 CFR 312.32 & Safety reporting & Document 13 §6, with the state clock in document 02 §6 \\
21 CFR 312.42 & Clinical holds & Document 01 §3(d), the first participant hold \\
21 CFR 312.57 & Sponsor recordkeeping for drug disposition & Document 13 §5 pharmacy \\
21 CFR 312.62 & Investigator recordkeeping and retention & Document 03 §5, at six years \\
21 CFR 812.140 & Device records & Document 06 §4 robot instance register \\
21 CFR 812.150 & Device reports & Document 02 §6 reportable robotic events \\
\end{xltabular}
\end{mvltable}
\tabcapl{tab:fdamap}{Fourteen requirements and the document in this package that
discharges each. Every row names a document rather than a policy, so a reviewer
can ask for the artifact by name.}

\section{ICH E6(R3) Mapping}

The good clinical practice spine is the adapted ICH E6(R3) framework published
in March 2026 rather than the source guideline, because the adaptation addresses
the physical artificial intelligence case the source does not
\cite{iche6r3adapt,iche6r3}. Where the two differ, the difference is stated.

\begin{mvtable}
\begin{tabularx}{\textwidth}{>{\raggedright\arraybackslash}p{3.4cm} >{\raggedright\arraybackslash}p{3.6cm} >{\raggedright\arraybackslash}X}
\headrow \thead{Principle} & \thead{Site implementation} & \thead{Where the adaptation differs from the source}\\
\midrule
Quality by design & Document 12 §6 & Adds the model version as a quality-critical factor \\
Investigator responsibility & Document 14 §2 & Names a delegated task set for a model governance role the source does not contemplate \\
Data governance & Document 03 §5 & Treats a completion as a source record \\
Risk-proportionate monitoring & Document 13 §6 & Adds envelope and interlock records to the monitoring scope \\
Records and reports & Document 03 §5 & Extends retention to six years \\
\end{tabularx}
\end{mvtable}
\tabcapl{tab:iche6map}{Five good clinical practice principles. Each row states
where the adapted framework goes beyond the source, because a reviewer who is
told only that a site follows good clinical practice learns nothing about the
model.}

\section{The Advisory-Only Boundary and Software Change Control}

This is the load-bearing section of the document. The boundary is written as a
rule that can be tested rather than as a commitment.

\begin{enumerate}
\item \textbf{What the model may produce.} A summary of source documents already
  in the record; a proposed differential; a proposed order set; a proposed
  eligibility assessment; a draft narrative for an adverse event; and a
  literature retrieval with citations.

\item \textbf{What the model may not do.} It may not transmit any output to a
  participant. It may not write to the electronic data capture system. It may not
  set, change, or confirm a dose. It may not command, configure, or release any
  robot motion. These four prohibitions are enforced at the interface, not by
  policy: the model process has no write credential to the data capture system
  and no route to the robot control network, which is segmented under document
  08 §5.

\item \textbf{What a human must do first.} A licensed clinician named in the
  delegation log must record a disposition of accepted, modified, or rejected
  before any completion is acted on. A completion with no recorded disposition
  is treated as rejected.

\item \textbf{What record proves it.} The audit trail entry under document 03
  §5, carrying the model version, the prompt, the completion, the clinician
  name, the disposition, and the timestamp to the second.
\end{enumerate}

Verification precedes generation: the site does not accept a completion into the
record on the strength of the model's confidence, and the verification-first
position is the one established in the 2026 legislative and assurance work
\cite{vvuqbill,cliniciantrust}.

Software change control follows a predetermined change control plan
\cite{fdasamd}. A change to weights, system prompt, retrieval corpus, decoding
parameters, or guardrail configuration is a new model version and requires a new
register entry under document 03 §3 before clinical use. The register is a single
artifact serving both the state obligation and the federal plan, which is stated
in document 03 §3 and is repeated here so that neither reader is surprised.

\section{Pre-Submission Sequence and Evidence Package}

\begin{mvtable}
\begin{tabularx}{\textwidth}{>{\raggedright\arraybackslash}p{0.9cm} >{\raggedright\arraybackslash}p{3.8cm} >{\raggedright\arraybackslash}X}
\headrow \thead{Step} & \thead{Action} & \thead{Evidence carried, and the limitation stated with it}\\
\midrule
1 & Pre-submission request \cite{fdapresub} & The protocol, the boundary rule in §5, and the segmentation drawing \\
2 & Credibility package & ASME V\&V 40 and ICH M15 aligned score of 81.9 across 55 verification tests \cite{daraxsims,asmevv40,ichm15}. Limitation: no patient-specific pharmacokinetic, pharmacodynamic, or tumor growth parameters \\
3 & Simulation evidence & Ten-arm run at 12.8 months median overall survival against 5.4 \cite{qspmetpancre}. Limitation: assumes no acquired resistance and ideal pharmacodynamics \\
4 & External comparison & RASolute 302 at 13.2 against 6.6 \cite{rasolute302}. Limitation: metastatic and previously treated; silent on the resectable setting \\
5 & Model documentation & The version register, the guardrail statement, and the change control plan \cite{fdasamd,fdaaiguidance2025} \\
6 & Investigational new drug submission & The filed application \cite{indapp} \\
\end{tabularx}
\end{mvtable}
\tabcapl{tab:presub}{The six-step sequence. Every evidence row carries the
limitation its own source stated, in the same cell as the result, because a
reviewer who has to look elsewhere for the caveat has been given a claim rather
than evidence.}
